%! Trigonometric Equations and Inequalities — Unit 7, Lesson 7.3 — Group Activity (Answer Key)
\documentclass[10pt]{article}
\usepackage{precalculus-article}
\usepackage{precalculus-key}   % pulls in -boxes; do NOT also load -boxes
\newcommand{\TallMath}[1]{$\displaystyle #1\rule[-1.4em]{0pt}{3.2em}$}

\begin{document}

\pageheader{Unit 7, Lesson 7.3}{Group Activity --- Answer Key}
\namepartnerperiod

\vspace{0.08in}

\begin{scenariobox}[Context: Solving Trigonometric Equations]{sky}
Trigonometric equations arise whenever a periodic model is set equal to
a specific output value. In this activity your group will solve trig
equations and inequalities in several settings. Choose \textbf{one tier}
below based on your readiness. Work together --- everyone should be ready
to explain your group's reasoning.
\end{scenariobox}

\vspace{0.08in}

%----------------------------------------------------------------------
% Tier R Key
%----------------------------------------------------------------------
\begin{tcolorbox}[colback=white, colframe=black!40,
    title=\textbf{Tier R --- Remediate (Key)},
    fonttitle=\bfseries, arc=2mm, left=3mm, right=3mm, top=2mm, bottom=2mm]

The general solution of $\sin\theta = \dfrac{1}{2}$ is already provided:
\[
  \theta = \frac{\pi}{6} + 2\pi k \qquad \text{or} \qquad \theta = \frac{5\pi}{6} + 2\pi k,
  \quad k \in \mathbb{Z}.
\]

\textbf{1.}\quad Family 1 solutions:

$k=0$: $\theta = $ \ans{$\tfrac{\pi}{6}$} \quad
$k=1$: $\theta = $ \ans{$\tfrac{\pi}{6}+2\pi = \tfrac{13\pi}{6}$} \quad
$k=2$: $\theta = $ \ans{$\tfrac{\pi}{6}+4\pi = \tfrac{25\pi}{6}$}

\vspace{0.08in}
\textbf{2.}\quad Family 2 solutions:

$k=0$: $\theta = $ \ans{$\tfrac{5\pi}{6}$} \quad
$k=1$: $\theta = $ \ans{$\tfrac{5\pi}{6}+2\pi = \tfrac{17\pi}{6}$} \quad
$k=2$: $\theta = $ \ans{$\tfrac{5\pi}{6}+4\pi = \tfrac{29\pi}{6}$}

\vspace{0.08in}
\textbf{3.}\quad Solutions in $[0, 4\pi]$: \ans{$\tfrac{\pi}{6},\,\tfrac{5\pi}{6},\,\tfrac{13\pi}{6},\,\tfrac{17\pi}{6}$} (4 solutions, since $4\pi \approx 12.57$ and $\tfrac{25\pi}{6} \approx 13.09 > 4\pi$).

How many? \ans{4}

\vspace{0.08in}
\textbf{4.}\quad $\theta = \dfrac{\pi}{2} + \pi k$ in $[0,2\pi]$:
\ans{$\theta = \dfrac{\pi}{2}$ and $\theta = \dfrac{3\pi}{2}$}

\vspace{0.08in}
\textbf{5.}\quad \ans{Trig functions are periodic, so the same output value repeats once per period. In $[0,4\pi]$ there are two full periods, giving two occurrences of each solution family.}

\end{tcolorbox}

\vspace{0.08in}

%----------------------------------------------------------------------
% Tier A Key
%----------------------------------------------------------------------
\begin{tcolorbox}[colback=white, colframe=black!40,
    title=\textbf{Tier A --- Apply (Key)},
    fonttitle=\bfseries, arc=2mm, left=3mm, right=3mm, top=2mm, bottom=2mm]

\textbf{1.}\quad $2\sin\theta + \sqrt{3} = 0$:

Step 1: \ans{$\sin\theta = -\dfrac{\sqrt{3}}{2}$}

Step 2: \ans{Principal: $\arcsin\!\bigl(-\tfrac{\sqrt{3}}{2}\bigr) = -\tfrac{\pi}{3}$ (Q4 reference $= \tfrac{\pi}{3}$). In $[0,2\pi)$: $\theta = \pi + \tfrac{\pi}{3} = \tfrac{4\pi}{3}$ (Q3) and $\theta = 2\pi - \tfrac{\pi}{3} = \tfrac{5\pi}{3}$ (Q4).}

Step 3: \ans{Sine is negative in Q3 and Q4.}

Step 4: \ans{$\theta = \dfrac{4\pi}{3}$ and $\theta = \dfrac{5\pi}{3}$}

\vspace{0.08in}
\textbf{2.}\quad $2\cos^2\!\theta - \cos\theta - 1 = 0$:

Factor: \ans{$(2\cos\theta + 1)(\cos\theta - 1) = 0$}

$u = \cos\theta$: \ans{$u = -\tfrac{1}{2}$} or \ans{$u = 1$}

For $\cos\theta = -\tfrac{1}{2}$: \ans{$\theta = \tfrac{2\pi}{3}$ and $\theta = \tfrac{4\pi}{3}$}

For $\cos\theta = 1$: \ans{$\theta = 0$}

All solutions in $[0,2\pi)$: \ans{$\theta = 0,\; \dfrac{2\pi}{3},\; \dfrac{4\pi}{3}$}

\vspace{0.08in}
\textbf{3.}\quad $d(t) = 2\cos\!\bigl(\tfrac{\pi}{6}t\bigr) + 4 = 3$:

(a) \ans{$\cos\!\bigl(\tfrac{\pi}{6}t\bigr) = -\tfrac{1}{2}$. Let $u = \tfrac{\pi}{6}t$: $u = \tfrac{2\pi}{3}$ or $u = \tfrac{4\pi}{3}$. Then $t = \dfrac{6}{\pi}\cdot\tfrac{2\pi}{3} = 4$ and $t = \dfrac{6}{\pi}\cdot\tfrac{4\pi}{3} = 8$. Solutions: $t = 4$ and $t = 8$ hours.}

(b) \ans{The general solution gives $t = 4 + 12k$ and $t = 8 + 12k$. For $k = 1$: $t = 16$ and $t = 20$, both outside $[0,12]$. Context restricts $t$ to one 12-hour period, so only $t = 4$ and $t = 8$ are valid.}

\end{tcolorbox}

\vspace{0.08in}

%----------------------------------------------------------------------
% Tier E Key
%----------------------------------------------------------------------
\begin{tcolorbox}[colback=white, colframe=black!40,
    title=\textbf{Tier E --- Extend (Key)},
    fonttitle=\bfseries, arc=2mm, left=3mm, right=3mm, top=2mm, bottom=2mm]

\textbf{1.}\quad $f(\theta) = 0$ at \ans{$\theta = \dfrac{\pi}{3}$ and $\theta = \dfrac{5\pi}{3}$}

\textbf{2.}\quad $f(\theta) > 0$ on \ans{$\bigl[0, \tfrac{\pi}{3}\bigr) \cup \bigl(\tfrac{5\pi}{3}, 2\pi\bigr]$} in $[0,2\pi]$.

\begin{teachernote}
From the table: $f(0)=1>0$, $f(\pi/3)=0$, $f(\pi/2)=-1<0$, $f(\pi)=-3<0$, $f(3\pi/2)=-1<0$, $f(5\pi/3)=0$, $f(2\pi)=1>0$.
Positive only near $0$ and $2\pi$.
\end{teachernote}

\textbf{3.}\quad $2\cos\theta - 1 > 0 \Rightarrow \cos\theta > \tfrac{1}{2}$.

\ans{Cosine $> \tfrac{1}{2}$ in $\bigl(-\tfrac{\pi}{3}, \tfrac{\pi}{3}\bigr)$ per period. In $[0,2\pi]$: $\cos\theta > \tfrac{1}{2}$ on $\bigl[0, \tfrac{\pi}{3}\bigr) \cup \bigl(\tfrac{5\pi}{3}, 2\pi\bigr]$. This confirms the table-based answer.}

\textbf{4a.}\quad $L(t) \ge 13.5$:

\ans{$3\sin\!\bigl(\tfrac{2\pi}{365}(t-80)\bigr) + 12 \ge 13.5 \Rightarrow \sin\!\bigl(\tfrac{2\pi}{365}(t-80)\bigr) \ge \tfrac{1}{2}$.\\
Let $u = \tfrac{2\pi}{365}(t-80)$. Boundary: $u = \tfrac{\pi}{6}$ and $u = \tfrac{5\pi}{6}$.\\
$t = 80 + \tfrac{365}{2\pi}\cdot u$.\\
$t_1 = 80 + \tfrac{365}{2\pi}\cdot\tfrac{\pi}{6} = 80 + \tfrac{365}{12} \approx 80 + 30 = 110$ (day 110).\\
$t_2 = 80 + \tfrac{365}{2\pi}\cdot\tfrac{5\pi}{6} = 80 + \tfrac{365 \cdot 5}{12} \approx 80 + 152 = 232$ (day 232).\\
$L(t) \ge 13.5$ for approximately days $110 \le t \le 232$.}

\textbf{4b.}\quad \ans{This represents roughly late April through mid-August, when the city has more than 13.5 hours of daylight. The general solution produces values like $t = 110 + 365k$ for other years ($k \ne 0$), but since $t \in [1,365]$ represents a single year, only $k=0$ is contextually valid.}

\end{tcolorbox}

\end{document}
